\section{Minimalisasi DFA}
\label{sec:ch06-minimasi}

\subsection{State yang Ekuivalen}

\begin{definisi}[Ekuivalensi state]
Dua state \(p\) dan \(q\) ekuivalen jika untuk setiap string \(w\),
\(\hat{\delta}(p,w) \in F\) tepat ketika
\(\hat{\delta}(q,w) \in F\). State yang tidak ekuivalen disebut
terbedakan.
\end{definisi}

Minimalisasi DFA berarti menggabungkan state-state ekuivalen sehingga
diperoleh DFA dengan jumlah state paling sedikit yang mengenali bahasa
yang sama. DFA minimal bersifat unik untuk suatu bahasa regular
(hingga penamaan state).

\subsection{Penyempurnaan Partisi}

Algoritma yang umum adalah penyempurnaan partisi:
\begin{enumerate}[label=\arabic*.,leftmargin=*]
  \item partisi awal: \(\{F, Q \setminus F\}\);
  \item untuk setiap grup, kelompokkan state berdasarkan vektor
        tujuan transisinya terhadap grup saat ini;
  \item ulangi sampai tidak ada grup yang terpecah;
  \item gabungkan state dalam grup yang sama.
\end{enumerate}

\begin{contoh}[Minimalisasi DFA \(abb\)]
Untuk DFA pada \figref{fig:ch06-dfa-abb}, partisi awal adalah
\(\{A,B,C\}\) dan \(\{D\}\). Pada iterasi pertama:
\begin{itemize}[leftmargin=*]
  \item \(A\) dan \(B\) memiliki transisi \(b\) ke grup
        \(\{A,B,C\}\), sedangkan \(C\) bertransisi \(b\) ke \(D\).
        Jadi \(C\) terpisah.
  \item \(A\) dan \(B\) dipisahkan karena \(A \xrightarrow{b} A\)
        (grup \(\{A,B\}\)) sedangkan \(B \xrightarrow{b} C\)
        (grup \(\{C\}\)).
\end{itemize}
Hasil akhirnya adalah empat grup \(\{A\}, \{B\}, \{C\}, \{D\}\).
DFA tersebut sudah minimal.
\end{contoh}

\subsection{Artefak: Tabel Minimalisasi}

Tabel \tabref{tab:ch06-minimasi} mencatat setiap iterasi penyempurnaan
partisi. Tabel ini adalah artefak Sub-CPMK 2.3; setiap pemisahan
harus disertai alasan berdasarkan transisi state.

\begin{table}[H]
\centering
\caption{Tabel minimalisasi DFA \((a \mid b)^{*}abb\)}
\label{tab:ch06-minimasi}
\begin{tabularx}{\textwidth}{@{}c l X X@{}}
\toprule
\textbf{Iterasi} & \textbf{Partisi} & \textbf{State yang dipisah} &
\textbf{Alasan} \\
\midrule
0 & \(\{A,B,C\}, \{D\}\) & --- & State penerima dipisahkan dari
state bukan penerima \\
1 & \(\{A,B\}, \{C\}, \{D\}\) & \(C\) & \(C \xrightarrow{b} D\),
sedangkan \(A,B \xrightarrow{b}\) ke \(\{A,B,C\}\) \\
2 & \(\{A\}, \{B\}, \{C\}, \{D\}\) & \(A\) dan \(B\) &
\(A \xrightarrow{b} A\), sedangkan \(B \xrightarrow{b} C\) \\
\bottomrule
\end{tabularx}
\end{table}

\begin{figure}[H]
\centering
\begin{tikzpicture}[node distance=6mm and 8mm]
  \node[phase, text width=2.4cm] (p0) {\(\{A,B,C\}\)};
  \node[phase, text width=2.0cm, right=of p0] (p1) {\(\{D\}\)};
  \node[phase, text width=1.8cm, below=of p0] (p2) {\(\{A,B\}\)};
  \node[phase, text width=1.8cm, right=of p2] (p3) {\(\{C\}\)};
  \node[phase, text width=1.8cm, below=of p2] (p4) {\(\{A\}\)};
  \node[phase, text width=1.8cm, right=of p4] (p5) {\(\{B\}\)};
  \draw[arrow] (p0) -- (p1);
  \draw[arrow] (p0) -- (p2);
  \draw[arrow] (p2) -- (p4);
  \draw[arrow] (p2) -- (p5);
\end{tikzpicture}
\caption{Penyempurnaan partisi pada DFA \((a \mid b)^{*}abb\)}
\label{fig:ch06-partisi}
\end{figure}

\subsection{Kompleksitas}

Implementasi sederhana penyempurnaan partisi memiliki kompleksitas
\(O(|Q|^2 |\Sigma|)\) untuk DFA dengan \(Q\) state. Algoritma
Hopcroft menurunkan kompleksitas menjadi
\(O(|Q| \log |Q| \cdot |\Sigma|)\). Untuk lexer praktis, ukuran DFA
biasanya kecil, sehingga implementasi sederhana sudah memadai.

\subsection{Implementasi}

Fungsi \code{minimize} pada \file{fa\_tool.py} melengkapi DFA dengan
sink state, membuang state tidak terjangkau, lalu menyempurnakan
partisi. Potongan awalnya ditunjukkan pada \lstref{lst:ch06-minimize}.

\includecode[
  language=Python,
  firstline=36,
  lastline=52,
  firstnumber=36,
  caption={Awal fungsi \code{minimize} pada \file{fa\_tool.py}},
  label={lst:ch06-minimize}
]{../../code/tools/fa_tool.py}

Sink state diperlukan agar setiap state memiliki transisi untuk semua
simbol. Tanpa sink state, dua state yang sama-sama menolak suatu
simbol dapat dianggap berbeda secara keliru.
